\chapter{Writing scripts}
\label{ch:scripts}

\section{What a script is}

A script is a Python file. It defines one function, \code{apply}, with one
parameter, the \code{vet} object. hdl-vet calls \code{apply(vet)} once.

\lstinputlisting[style=script]{examples/scripts/first.py}

hdl-vet reads the script as text and runs it. It does not import it, so it
leaves no \file{__pycache__} folder next to it. A script that does not define
\code{apply} is an error. So is a script that you name twice.

Whatever a script prints goes to the standard output of \cmd{vet}. Section~\ref{sec:program}
says what this means for your safety.

A script can fail in three ways. All of them exit with status 2 and write nothing
(chapter~\ref{ch:cli} lists the exit statuses).

\begin{itemize}
\item An operation refuses, or finds drift. hdl-vet prints \code{vet: script NAME:} and the
  reason. The error is a \code{VetError}, or a \code{DriftError}, which is a kind of
  \code{VetError}.
\item The script raises another Python exception. hdl-vet prints
  \code{vet: script NAME raised TYPE: MESSAGE} and then a Python traceback.
\item A raw edit with \apiname{Vet.splice} overlaps an earlier one. The error is a
  \code{SpliceError}. It is not a \code{VetError}, so \code{except VetError} in your script
  does not catch it, and the message is only \code{vet: [2, 7) overlaps [0, 5)}, with the
  two ranges and without the name of the script.
\end{itemize}

\section{Find modules and instances}

You start with \sig{Vet.module}. It takes the name of a module and gives you a
\term{module handle}. You edit a module through its handle. From a module handle
you get \term{instance handles} with \sig{Module.instance} (one instance, by name)
and \sig{Module.instances} (all instances that match a filter).

\lstinputlisting[style=script]{examples/scripts/find.py}

The next two scripts fail on purpose. The first has a spelling mistake in a module
name. The second edits \code{timer} without \code{instances}. In the transcript, look at
the output of \file{find.py}: it prints the paths and the names that the list below
explains. Then look at the two error messages.

\lstinputlisting[style=script]{examples/scripts/typo.py}
\lstinputlisting[style=script]{examples/scripts/needs_count.py}

\lstinputlisting[style=console]{examples/console/find.txt}

\begin{itemize}
\item \apiname{Vet.module} refuses a name that the design does not have, and suggests
  the closest names. It also refuses a module that is not instantiated under the
  top module, and a module that comes from an added file
  (\apiname{Vet.add_file}, chapter~\ref{ch:operations}). hdl-vet edits only the
  modules that it can see in the elaborated design.
\item \apiname{Module.instance} takes the name of one instance in the module. It
  refuses a name that is not there, with suggestions.
\item \apiname{Module.instances} takes an optional \code{module} (the module that is
  instantiated, an exact name), an optional \code{name} (a pattern with
  \code{*} and \code{?}), and \code{count}, which is required: you say how many
  instances you expect. A different number is an error that lists what matched.
  This protects you when a generator changes the number of instances.
\item \sig{Module.paths} is the list of the paths of the module in the
  design, for example \code{soc_top.u_periph.u_timer0}. The first part is the
  name of the top module. \sig{Instance.paths} lists the paths of one instance, one
  for each path of its module. \sig{Vet.module_paths} gives the paths of any
  module by name. These paths describe the design as hdl-vet read it at the start
  of the run. They do not change when your script edits the design.
\item \code{top.path} is the \emph{file} that holds the module. It is not one of the
  paths above. It is an attribute of the model that \code{top.info} returns, and the
  handle forwards it (see the last item). Scripts use it to read and edit the text of the
  module with \apiname{Vet.text} and \apiname{Vet.splice}.
\item \sig{Module.name} and \sig{Instance.name} are the names. \sig{Instance.target}
  is \code{module.instance}, the text that hdl-vet uses in its log and its
  messages.
\item \sig{Module.info} and \sig{Instance.info} give the model of the compiler for
  the module or the instance: ports, instances, connections, and so on. The
  handles forward attribute lookups to it, so \code{top.info.ports} and
  \code{top.ports} are the same list. Read these models; do not change them.
  A handle looks the design up again at each call, so it is never out of date.
\end{itemize}

\subsection*{A module that exists more than once}

\code{timer} is instantiated twice, so \code{periph_wrap} has two instances of
it, and the design has two copies of \code{timer}. When you edit a module, the edit
changes every copy. hdl-vet makes you say that you know.

\begin{itemize}
\item \code{vet.module("timer", instances=2)} confirms that you expect exactly
  two copies. Without \code{instances}, an edit of \code{timer} is refused, with a
  message that lists the paths (the last command of the transcript above). If the
  number is wrong, \apiname{Vet.module} refuses at once.
\item The same holds for instance operations. An instance inside \code{timer}
  would be inside two copies of \code{timer}. You get the handle through
  \code{vet.module("timer", instances=2)}.
\end{itemize}

\section{Applied and skipped}

Every operation looks at the design before it edits.

\begin{enumerate}
\item If the design already shows the result that the operation asks for, the
  operation is \term{skipped}. It queues nothing, and it registers no
  expectations.
\item If the design is in the state that the operation expects before the edit, the
  operation is \term{applied}: it queues the edit and registers what it
  expects of the sign-off.
\item Otherwise there is \term{drift}: the operation stops with an error that
  says what it found.
\end{enumerate}

Because of the first rule you can run a script a second time on a design that
the script has already changed. Nothing happens twice. The summary line,
\code{N applied, M skipped}, counts the operations of the run. A raw text edit
with \apiname{Vet.splice} is not an operation and is not counted.

\lstinputlisting[style=script]{examples/scripts/skipped.py}

In the next transcript, \file{skipped.py} runs after \file{first.py}. It asks
for the port that \file{first.py} has added, so that operation is skipped. The
summary is \code{3 applied, 1 skipped}.

\lstinputlisting[style=console]{examples/console/applied_skipped.txt}

\subsection*{The argument \texttt{was}}

An operation that \emph{replaces} something takes an optional argument \code{was}.
It says what you expect to find. The operations are \apiname{Instance.connect},
\apiname{Instance.set_param}, \apiname{Instance.disconnect}, \apiname{Module.remove_assign}
and \apiname{Instance.replace_module}. \apiname{Module.remove_instance} has the argument
\code{module}, which works in the same way as \code{was}.

\begin{itemize}
\item If the design shows the new value already, the operation is skipped.
  \code{was} is not looked at.
\item If the design shows exactly the value in \code{was}, the operation
  replaces it.
\item If the design shows another value, and \code{was} is given but does not match,
  the operation stops with a drift error. The error quotes the value that it found.
  This holds for all of these operations.
\item If the design shows another value and \code{was} is \emph{missing}, the
  operations differ. \apiname{Instance.connect} and \apiname{Instance.set_param} stop
  with a drift error, because they would overwrite a value that you did not name.
  \apiname{Instance.disconnect}, \apiname{Module.remove_assign},
  \apiname{Instance.replace_module} and \apiname{Module.remove_instance} do not
  check: without \code{was} (or \code{module}) they do their edit on whatever they
  find.
\end{itemize}

The values are compared as text, by the tokens of the language. White space and
comments do not count. You may write a number: \code{was=2} is the text
\code{2}.

The last command of the transcript runs this script. The design does not show
\code{1'b1} at \code{test_i}, so the operation stops with exit status 2, and
nothing is written.

\lstinputlisting[style=script]{examples/scripts/drift.py}

Drift is an error of the script, not a failed sign-off. The exit status is 2.
Give \code{was} always. It turns a silent wrong edit into a message.

\section{Several scripts}

You give scripts to \cmd{vet apply} with \opt{-s}, as many as you like. If the
project file lists scripts, they run first, and the scripts of the command line
run after them. The scripts run in order, in one session. A script sees what the
earlier scripts have done: when a script has finished, hdl-vet loads the edited
text, so the next script finds the new ports, wires and connections. The sign-off
runs once, at the end, over all scripts.

\section{Groups}

Some operations check themselves. They are the \term{checked operations}:
\apiname{Vet.add_file}, \apiname{Instance.replace_module}, \apiname{Module.remove_instance},
\apiname{Module.remove_wire}, \apiname{Module.remove_assign}, \apiname{Module.remove_port} and
\apiname{Module.remove_param}. A checked operation first applies and loads everything
that was queued. Then it makes its edit and loads the design again. It is
\emph{refused}, and everything it did is taken back, when the edited design

\begin{itemize}
\item cannot be loaded, or
\item has a new diagnostic (a new error or a new warning), or
\item leaves a signal or a port \term{dangling}: something reads it, but nothing
  drives it, or
\item leaves a signal partly undriven, for a removal that took away a driver.
\end{itemize}

A refusal is an error of the script (exit status 2). The message starts with the
operation and its target. After a refusal the session is as it was before the
call. If the script catches the error, it can go on.

A removal often needs a second edit to be clean: you remove the driver of a
signal, and you must also remove the signal or its reader. Each step alone is
refused. A \term{group} makes the comparison once, at the end of the block. You open
it with \code{with vet.group():}, which is \sig{Vet.group}. If the comparison fails, or if
an exception leaves the block, everything done in the block is taken back. A group inside
a group is part of the outer group.

\lstinputlisting[style=script]{examples/scripts/no_group.py}
\lstinputlisting[style=script]{examples/scripts/group.py}

The first script removes the reset synchroniser, and does not reconnect
\code{u_periph}. The wire \code{rst_sync_n} is then read but not driven: the
first operation is refused. The second script does all three steps in a group,
and succeeds.

\lstinputlisting[style=console]{examples/console/group.txt}

\subsection*{Notes}

When a checked operation, or a group, leaves something that nobody uses, it is
not an error. hdl-vet adds a line to the list \code{vet.notes}, and the command
prints it as \code{note: ...}. A note starts with the operation and its target,
or with \code{group:} for a group. Notes never make a run fail. The transcript of
\code{retire.py} in chapter~\ref{ch:operations} shows one.

\section{Scopes}

A \term{scope} limits part of a script to one module and what is below it. It
protects you from a mistake that edits another part of the design. You
open a scope with \sig{Vet.scope}, as a \code{with} block.

\lstinputlisting[style=script]{examples/scripts/scope_periph.py}

The scope is the module \code{periph_wrap}. Everything instantiated below it,
here the two timers, is in the scope too. The design is still the whole design,
and the sign-off still judges all of it. A scope adds refusals. It takes no check
away.

\subsection*{What is inside}

\begin{itemize}
\item On entry, hdl-vet applies and loads everything that is queued, as for a
  group. The module must be one that \apiname{Vet.module} accepts, and it must be
  instantiated under the top module. The top module itself is allowed. The scope
  then holds the whole design.
\item An \emph{instance} is inside when it is an instance of the scope module, or
  lies below one, in the design as it is when you enter the block.
\item A \emph{module} is inside when it has at least one instance, and every
  instance of it is inside. A module that is also instantiated elsewhere is not
  inside, even if some of its instances are.
\end{itemize}

In the example, \code{periph_wrap} and \code{timer} are inside the scope
\code{periph_wrap}. \code{soc_top} and \code{rst_sync} are not.

\subsection*{What is refused}

\begin{itemize}
\item An operation of a module handle or an instance handle that would edit the
  text of a module that is not inside is refused when you call it. This holds for
  a handle that you got before the block, too. An instance operation edits the
  module in which the instance is written.
\item \apiname{Module.remove_port} and \apiname{Module.remove_param} also edit every
  module that instantiates the module. They are refused when one of those modules is
  not inside. For the scope module this is the case, unless it is the top module.
\item A refusal names the operation, the module, the scope, and where the module is
  instantiated. The scope is checked before anything else, so an operation that would
  have been skipped is refused too. Nothing is queued. The block goes on, if your script
  catches the error.
\item Reading is free. \apiname{Vet.module}, \apiname{Module.info}, \apiname{Module.paths} and
  \apiname{Module.instance} work for every module.
\item \apiname{Vet.splice} in a scope is refused unless the whole range lies in
  the text of one module that is inside, from the keyword \code{module} to the
  end of \code{endmodule}. It is refused too when the new text holds one of the
  compiler directives \code{`include}, \code{`define}, \code{`undef},
  \code{`undefineall}, \code{`celldefine}, \code{`endcelldefine}, \code{`line},
  \code{`default_nettype}, \code{`unconnected_drive},
  \code{`nounconnected_drive}, \code{`timescale}, \code{`resetall},
  \code{`pragma}, \code{`begin_keywords} or \code{`end_keywords}, outside a
  comment and outside a string. These directives reach text outside the
  module. A use of a macro and the directives \code{`ifdef}, \code{`ifndef},
  \code{`elsif}, \code{`else} and \code{`endif} are allowed.
\item \apiname{Vet.add_file}, \apiname{Vet.expect} and \apiname{Vet.permit} are not
  restricted.
\item \apiname{Vet.punch_port} and \apiname{Vet.propagate_param} start at the scope module.
  The path in \code{to} begins with the name of the scope module, and the new port
  or parameter ends in the list of the scope module itself. Nothing above it is
  edited. In the example, \code{to="periph_wrap.u_timer0"}. A path that starts
  anywhere else is refused, and the scope module alone is refused. Both operations
  are refused when the scope module is instantiated more than once.
\item A scope cannot be opened inside a scope or inside a group. A group may be
  opened inside a scope. \apiname{Vet.check}, \apiname{Vet.save} and \apiname{Vet.finish}
  are refused inside a scope.
\end{itemize}

The scope module's own ports and parameters may change. The parent's
instantiation is outside the scope. So a new input port is not connected in the
parent, and the sign-off fails, until a later operation outside the block
connects it. The example does that after the block: it adds the port to the top
module and connects \code{u_periph}.

\subsection*{What is checked at the end}

When the block ends, hdl-vet loads the design again and compares it with the
design on entry. It raises an error that says \emph{outside the scope}, and takes
back everything done in the block, when

\begin{itemize}
\item any byte outside the text of the inside modules has changed, or
\item the structure has changed at an instance that is not inside, or
\item a module that is inside names, by a hierarchical path, a signal, a task or a
  function of an instance outside the scope, and the design on entry had no such
  reference in that module.
\end{itemize}

The message names one such reference, even if the block added several. The compiler
does not look at a name in code that it does not elaborate, for example an unused
sequence or property, so such a name is not seen. The last point holds however the text
got there: by a splice, or in the
expression of \apiname{Instance.connect} or \apiname{Module.add_assign}. A path that ends
at or below an instance of the scope module is fine. A name that goes through an
interface port of the module is a use of the port, not a reach outside.
If an exception leaves the block, everything done in it is taken back, as for a
group. A script must close every block it opens: if \code{apply} returns with a
scope or a group still open, \cmd{vet} stops with exit status 2 and names the
script.

The rest of the design can be edited again after the block. A scope is not kept in
a saved session (chapter~\ref{ch:session}). Operations in a scope are in
\code{vet.log} as they are outside of one. The next scripts show a second scope,
and two ways to leave one.

\lstinputlisting[style=script]{examples/scripts/scope_reset.py}
\lstinputlisting[style=script]{examples/scripts/scope_leaves.py}
\lstinputlisting[style=script]{examples/scripts/scope_reaches.py}

\lstinputlisting[style=console]{examples/console/scope.txt}

In the transcript, the first command runs two scripts, each with its own scope. It adds a parameter
to \code{rst_sync} inside a scope, and overrides it on the instance outside the
block. The second script writes to \code{soc_top} inside the scope
\code{periph_wrap}, and is refused at the call. The third script names an
instance of \code{rst_sync} from inside \code{periph_wrap}. Without the scope,
the sign-off accepts this reference. In the scope, it reaches outside, so the
block is refused and taken back.

\section{A script is a program}
\label{sec:program}

A script is a Python program. It runs with your rights: it can read and write any file
that you can. hdl-vet does not run it in a sandbox. The sign-off checks the design that
comes out, not what the script did on your machine. \cmd{vet apply} runs the scripts also
with \opt{--dry-run}, and a project file lists scripts that run when you type
\cmd{vet apply}. Run only scripts and projects that you trust.
